% ======================================================================
% Appendix C -- Templates
% ======================================================================

\chapter{Templates}
\label{app:templates}

These templates collect the deliverables used in every case. Editable versions---for documents, slides and the TikZ notation styles---are in the companion repository's \texttt{templates/} folder.

\newtcolorbox{template}[1]{enhanced, breakable, colback=white, colframe=black!50, boxrule=0.5pt, arc=1pt, title={#1}, fonttitle=\bfseries\small, fontupper=\small, coltitle=black, colbacktitle=codefill}

\section{Solution graph checklist}

\begin{template}{Solution graph: well-formedness (Chapter~\ref{ch:toolkit})}
\begin{enumerate}[label=\textbf{W\arabic*.}, leftmargin=2.3em, itemsep=1pt]
  \item Every agent states its autonomy level\index{autonomy spectrum} and stop condition.
  \item Every write action has a checkpoint or guardrail before it, or a note saying why not.
  \item No agent touches a system of record except through an MCP server.
  \item Thresholds and rules live in code nodes, not agent annotations.
  \item Every knowledge store names an owner and a freshness target.
  \item Every edge into an agent that crosses a trust boundary passes a guardrail.
  \item Every branch has its condition on the edge, and branches cover all cases.
  \item Every node (except triggers, users and systems of record) has a trace tag.
\end{enumerate}
\end{template}

\section{Node card}

\begin{template}{Node card: \emph{name}}
\begin{tabularx}{\linewidth}{@{}>{\bfseries}p{3.4cm}X@{}}
Purpose & One sentence. \\
Inputs & What it receives, and from where. \\
Outputs & Schema or format. \\
Autonomy and stop & Level 0--5; stop condition; limits on calls and time. \\
Tools & Each tool, read or write, and its server's owner. \\
Model tier & Smallest model that meets the quality target. \\
Quality target & Metric, threshold, evaluation set. \\
Cost and latency budget & Per call or per task; p90. \\
Failure handling & What happens on each kind of failure. \\
Traces to & KPIs and risks. \\
\end{tabularx}
\end{template}

\section{Trade-off table}

\begin{template}{Trade-off table}
\begin{tabular}{@{}lccc@{}}
Criterion & Weight & D1 & D2 \ldots \\
\hline
Money rows: run cost, human time, expected error cost (per unit) & --- & & \\
Quality & & & \\
Speed & & & \\
Reliability & & & \\
Security and privacy & & & \\
Control and auditability & & & \\
Maintainability & & & \\
Human workload (if not in money rows) & & & \\
\hline
Weighted score & 1.00 & & \\
\end{tabular}

\smallskip
Constraints checked first: \underline{\hspace{8cm}}\\
Sensitivity: share of perturbed weightings each design wins; what would flip the decision.
\end{template}

\section{Failure walk-through}

\begin{template}{Failure walk-through}
\begin{tabular}{@{}llllll@{}}
Node & Failure & Effect & Detection & Mitigation & S\,/\,L \\
\hline
 & & & & & \\
\end{tabular}

\smallskip
Include failures of tools, data and people, not only models. Mark the rows that changed the design.
\end{template}

\section{Decision record}

\begin{template}{DR-\emph{area}-\emph{number}: \emph{decision in one line}}
\textbf{Status and date.} Proposed / accepted / superseded; owner.\\
\textbf{Context.} The problem, KPIs, risks and constraints, in a few lines.\\
\textbf{Options considered.} Each with a solution graph reference.\\
\textbf{Decision.}\\
\textbf{Reasons.} Including the sensitivity result.\\
\textbf{Consequences.} What changes for whom; what new obligations it creates.\\
\textbf{Reversal triggers\index{reversal trigger}.} Observable conditions under which this decision is reopened.
\end{template}

\section{One-page memo}

\begin{template}{Memo to the executive sponsor}
\textbf{To / From / Re / Decision needed by.}\\
\textbf{Recommendation.} What you ask for, and its cost.\\
\textbf{The problem.} In the sponsor's numbers, from the process map\index{current-process map}.\\
\textbf{What changes.} In business terms; no framework or model names.\\
\textbf{Value.} Floor case\index{floor case} from measured inputs; upside to be tested; the assumption the case rests on.\\
\textbf{Risks and safeguards.} The two or three that matter, and how they are controlled.\\
\textbf{Pilot and decision.} Scope, comparison, duration, go/no-go\index{go/no-go criteria} criteria, owner.\\
\textbf{What we need from you.} Including any decision only the sponsor can make.
\end{template}

\section{Go/no-go criteria}

\begin{template}{Pilot go/no-go}
\begin{tabular}{@{}lll@{}}
Criterion & Baseline & Go if \\
\hline
Each KPI & & \\
Each error budget (zero for catastrophic errors; also a stop rule) & & \\
Adoption & & \\
\end{tabular}

\smallskip
Stop rules\index{stop rule}: \underline{\hspace{9cm}}\\
Owner of the decision: \underline{\hspace{8cm}}
\end{template}
